% Figure: the 22 municipalities, shaded by home codes per square kilometre.
% Data: data/map.dat, data/municipalities.dat, data/islands.dat (tools/paper_figures.py).
\documentclass[tikz,border=2mm]{standalone}
% Shared preamble of the figures: fonts and the paper's colours.
\usepackage{fontspec}
\setmainfont{LibertinusSans}[Extension=.otf,UprightFont=*-Regular,BoldFont=*-Bold,ItalicFont=*-Italic]
\usepackage{pgfplots,pgfplotstable}
\pgfplotsset{compat=1.18}
\usetikzlibrary{calc,patterns}
\definecolor{ink}{HTML}{003893}
\definecolor{flagred}{HTML}{CF2027}
\definecolor{star}{HTML}{F7D116}
\definecolor{sea}{HTML}{EEF3FA}
\definecolor{grid}{HTML}{9AA7BD}
\pgfplotsset{every axis/.append style={font=\small,tick label style={font=\footnotesize},label style={font=\small}}}
% Language of the labels: a figure compiled through a wrapper such as map-pt.tex
% (\def\figlang{pt}% Figure: the 22 municipalities, shaded by home codes per square kilometre.
% Data: data/map.dat, data/municipalities.dat, data/islands.dat (tools/paper_figures.py).
\documentclass[tikz,border=2mm]{standalone}
% Shared preamble of the figures: fonts and the paper's colours.
\usepackage{fontspec}
\setmainfont{LibertinusSans}[Extension=.otf,UprightFont=*-Regular,BoldFont=*-Bold,ItalicFont=*-Italic]
\usepackage{pgfplots,pgfplotstable}
\pgfplotsset{compat=1.18}
\usetikzlibrary{calc,patterns}
\definecolor{ink}{HTML}{003893}
\definecolor{flagred}{HTML}{CF2027}
\definecolor{star}{HTML}{F7D116}
\definecolor{sea}{HTML}{EEF3FA}
\definecolor{grid}{HTML}{9AA7BD}
\pgfplotsset{every axis/.append style={font=\small,tick label style={font=\footnotesize},label style={font=\small}}}
% Language of the labels: a figure compiled through a wrapper such as map-pt.tex
% (\def\figlang{pt}\input{map.tex}) prints its Portuguese labels, map-fr.tex its French ones.
\providecommand{\figlang}{en}
\def\langpt{pt}\def\langfr{fr}
\newcommand{\lbl}[3]{\ifx\figlang\langpt #2\else\ifx\figlang\langfr #3\else #1\fi\fi}
\newcommand{\thousands}{\ifx\figlang\langpt labelsp\else\ifx\figlang\langfr labelsp\else label\fi\fi}

\begin{document}
\pgfplotstableread{data/municipalities.dat}\munis
\pgfplotstableread{data/islands.dat}\isles
\begin{tikzpicture}
\begin{axis}[axis equal image, hide axis, width=15cm, xmin=-150, xmax=150, ymin=-140, ymax=140, clip=false]
\pgfplotsinvokeforeach{0,...,21}{
  \pgfplotstablegetelem{#1}{shade}\of\munis
  \edef\cmd{\noexpand\addplot[fill=ink!\pgfplotsretval!white, draw=white, line width=0.35pt, unbounded coords=jump]
    table[x=x, y=y, restrict expr to domain={\noexpand\thisrow{id}}{#1:#1}] {data/map.dat};}
  \cmd
}
\pgfplotsinvokeforeach{0,...,8}{
  \pgfplotstablegetelem{#1}{island}\of\isles \let\nm\pgfplotsretval
  \pgfplotstablegetelem{#1}{\thousands}\of\isles \let\cd\pgfplotsretval
  \pgfplotstablegetelem{#1}{lx}\of\isles \let\lx\pgfplotsretval
  \pgfplotstablegetelem{#1}{ly}\of\isles \let\ly\pgfplotsretval
  \pgfplotstablegetelem{#1}{anchor}\of\isles \let\an\pgfplotsretval
  \edef\cmd{\noexpand\node[font=\noexpand\footnotesize, align=center, anchor=\an, inner sep=1pt]
    at (axis cs:\lx,\ly) {\noexpand\textbf{\nm}\\\cd~\lbl{codes}{códigos}{codes}};}
  \cmd
}
\draw[gray] (axis cs:-148,-8) -- (axis cs:-98,-8) node[midway, below, font=\scriptsize] {50 km};
\draw[gray] (axis cs:-148,-10) -- (axis cs:-148,-6) (axis cs:-98,-10) -- (axis cs:-98,-6);
\end{axis}
\begin{scope}[shift={(0.3,5.2)}]
  \node[anchor=west, font=\footnotesize] at (0,1.45) {\lbl{Home codes per km\textsuperscript{2}}{Códigos de casa por km\textsuperscript{2}}{Codes de maison par km\textsuperscript{2}}};
  \foreach \s/\l [count=\k from 0] in {20/8,40/25,60/70,80/180,92/320} {
    \fill[ink!\s!white] (\k*0.9,0.6) rectangle ++(0.85,0.4);
    \node[font=\scriptsize, anchor=north] at (\k*0.9+0.425,0.58) {\l};
  }
\end{scope}
\end{tikzpicture}
\end{document}
) prints its Portuguese labels, map-fr.tex its French ones.
\providecommand{\figlang}{en}
\def\langpt{pt}\def\langfr{fr}
\newcommand{\lbl}[3]{\ifx\figlang\langpt #2\else\ifx\figlang\langfr #3\else #1\fi\fi}
\newcommand{\thousands}{\ifx\figlang\langpt labelsp\else\ifx\figlang\langfr labelsp\else label\fi\fi}

\begin{document}
\pgfplotstableread{data/municipalities.dat}\munis
\pgfplotstableread{data/islands.dat}\isles
\begin{tikzpicture}
\begin{axis}[axis equal image, hide axis, width=15cm, xmin=-150, xmax=150, ymin=-140, ymax=140, clip=false]
\pgfplotsinvokeforeach{0,...,21}{
  \pgfplotstablegetelem{#1}{shade}\of\munis
  \edef\cmd{\noexpand\addplot[fill=ink!\pgfplotsretval!white, draw=white, line width=0.35pt, unbounded coords=jump]
    table[x=x, y=y, restrict expr to domain={\noexpand\thisrow{id}}{#1:#1}] {data/map.dat};}
  \cmd
}
\pgfplotsinvokeforeach{0,...,8}{
  \pgfplotstablegetelem{#1}{island}\of\isles \let\nm\pgfplotsretval
  \pgfplotstablegetelem{#1}{\thousands}\of\isles \let\cd\pgfplotsretval
  \pgfplotstablegetelem{#1}{lx}\of\isles \let\lx\pgfplotsretval
  \pgfplotstablegetelem{#1}{ly}\of\isles \let\ly\pgfplotsretval
  \pgfplotstablegetelem{#1}{anchor}\of\isles \let\an\pgfplotsretval
  \edef\cmd{\noexpand\node[font=\noexpand\footnotesize, align=center, anchor=\an, inner sep=1pt]
    at (axis cs:\lx,\ly) {\noexpand\textbf{\nm}\\\cd~\lbl{codes}{códigos}{codes}};}
  \cmd
}
\draw[gray] (axis cs:-148,-8) -- (axis cs:-98,-8) node[midway, below, font=\scriptsize] {50 km};
\draw[gray] (axis cs:-148,-10) -- (axis cs:-148,-6) (axis cs:-98,-10) -- (axis cs:-98,-6);
\end{axis}
\begin{scope}[shift={(0.3,5.2)}]
  \node[anchor=west, font=\footnotesize] at (0,1.45) {\lbl{Home codes per km\textsuperscript{2}}{Códigos de casa por km\textsuperscript{2}}{Codes de maison par km\textsuperscript{2}}};
  \foreach \s/\l [count=\k from 0] in {20/8,40/25,60/70,80/180,92/320} {
    \fill[ink!\s!white] (\k*0.9,0.6) rectangle ++(0.85,0.4);
    \node[font=\scriptsize, anchor=north] at (\k*0.9+0.425,0.58) {\l};
  }
\end{scope}
\end{tikzpicture}
\end{document}
